\section{Introduction}
\label{sec:intro}

A phone camera at night faces a hard trade-off. A long exposure collects enough light, but the camera must stay still for seconds and anything that moves blurs. A short exposure freezes motion, but at 1/30\,s in a dim room the sensor records only a few photons per pixel. Shot noise, read noise and quantisation then dominate the signal, colours shift, and dark regions lose their structure. Restoring a clean, correctly coloured image from such a frame matters for night photography, surveillance, driving, drones and search-and-rescue. It is also a hard inverse problem, because much of the scene information is simply missing from the input.

The field moved in three steps. Classical pipelines denoise and then tone-map, or merge a burst of frames~\cite{dabov2007bm3d,hasinoff2016burst}. Methods that work on the processed sRGB image~\cite{wei2018retinexnet,guo2020zerodce,jiang2021enlightengan} break down in \emph{extreme} low light, because the camera's image processor has already quantised and amplified the noise. Chen et al.~\cite{chen2018sid} showed that a single U-Net~\cite{ronneberger2015unet} trained on \emph{RAW} short exposures and their long-exposure counterparts (the SID dataset) beats these pipelines by a wide margin. Two problems remained. First, paired data are camera-specific: a SID-trained model degrades on other sensors, and capturing new pairs for every camera is expensive. ELD~\cite{wei2020eld} addresses this with calibrated noise synthesis, and Prabhakar et al.~\cite{prabhakar2021fewshot} with few-shot adaptation. Second, the regression network outputs a single blurry average when information is missing. Diffusion models~\cite{ho2020ddpm,saharia2022sr3}, and in particular \emph{diffusion bridges} that diffuse directly between the degraded and the clean image~\cite{liu2023i2sb,zhou2024ddbm}, promise sharper and more faithful outputs, at the cost of many network evaluations.

This project builds on the SID U-Net and asks three questions:
\begin{enumerate}
\item \textbf{Can the U-Net be told what it is looking at?} In extreme low light, noise and content mix most at the bottleneck. We add a \emph{bridge} module there that refines the deepest features before decoding. It is conditioned by FiLM~\cite{perez2018film} on the exposure ratio and an estimated noise level. During training it is pulled towards the features of the clean long exposure by teacher-latent distillation~\cite{romero2015fitnets}, which costs nothing at test time.
\item \textbf{Can it adapt to a new camera from a few scenes?} We compare full fine-tuning on ELD cameras with $K$-shot adaptation ($K=1,3,5$ scenes) from a supervised and a meta-learned (MAML~\cite{finn2017maml}) initialisation. We redesign the meta-learning tasks so that support and query never share a scene, and we state exactly how many scenes and samples each task uses.
\item \textbf{Does a diffusion bridge help?} We compare the feature-space ``DDBM-Feature'' model of the original project with a real pixel-space Denoising Diffusion Bridge Model~\cite{zhou2024ddbm} trained by bridge score matching.
\end{enumerate}

\paragraph{Takeaways.} \TAKEAWAYS

\paragraph{Changes from the first version.} The original code of this project was lost. Everything here was re-implemented from the first report, retrained and re-evaluated on full-resolution test images. In the process we found and fixed a flaw in the original meta-learning setup: support and query targets coincided. We also added the SID baseline, ablations and the pixel-space DDBM, which the first report announced but did not include.
